% Part: first-order-logic
% Chapter: models-theories
% Section: introduction

\documentclass[../../../include/open-logic-section]{subfiles}

\begin{document}

\olfileid[pt]{fol}{mat}{int}
\olsection{Introdução}

\begin{explain}
O desenvolvimento do método axiomático é uma conquista significativa
na história da ciência, e é de especial importância na história
da matemática.  Um desenvolvimento axiomático de um campo envolve o
esclarecimento de muitas questões: Do que trata o campo?  Quais são os
conceitos mais fundamentais?  Como eles se relacionam?  Todos os
conceitos do campo podem ser definidos em termos desses conceitos
fundamentais?  Que leis esses conceitos obedecem, e devem obedecer?

O método axiomático e a lógica foram feitos um para o outro.  A lógica formal
fornece as ferramentas para formular teorias axiomáticas, para provar
teoremas a partir dos axiomas da teoria de uma maneira precisamente especificada,
para estudar as propriedades de todos os sistemas que satisfazem os axiomas de uma
maneira sistemática.
\end{explain}

\begin{defn}
Um conjunto de !!{sentence}s~$\Gamma$ é \emph{fechado} sse, sempre que
$\Gamma \Entails !A$ então $!A \in \Gamma$.  O \emph{fecho} de um conjunto
de !!{sentence}s~$\Gamma$ é $\Setabs{!A}{\Gamma \Entails !A}$.

Dizemos que~$\Gamma$ é \emph{axiomatizado por} um conjunto de
sentenças~$\Delta$ se $\Gamma$ é o fecho de~$\Delta$.
\end{defn}

\begin{explain}
Podemos pensar em uma teoria axiomática como o conjunto de sentenças que é
axiomatizado por seu conjunto de axiomas~$\Delta$.  Em outras palavras, quando
temos uma linguagem de primeira ordem que contém símbolos não lógicos para os
primitivos da ciência desenvolvida axiomaticamente que desejamos estudar,
juntamente com um conjunto de !!{sentence}s que expressam as leis fundamentais
da ciência, podemos pensar na teoria como representada por todas as
!!{sentence}s nessa linguagem que são consequência dos axiomas.  Isso
varia de exemplos simples com apenas um único primitivo e
axiomas simples, como a teoria das ordens parciais, até teorias complexas
como a mecânica newtoniana.

Os fatos lógicos importantes que tornam essa abordagem formal do
método axiomático tão importante são os seguintes.  Suponha que $\Gamma$ é
um sistema de axiomas para uma teoria, isto é, um conjunto de sentenças.
\begin{enumerate}
\item Podemos afirmar precisamente quando um sistema de axiomas captura uma classe
  pretendida de !!{structure}s.  Isto é, se estamos interessados em uma certa
  classe de !!{structure}s, capturaremos com sucesso essa classe por
  um sistema de axiomas~$\Gamma$ sse as !!{structure}s são exatamente
  aquelas~$\Struct M$ tais que $\Sat{M}{\Gamma}$.
\item Podemos falhar nesse aspecto porque há $\Struct M$ tais
  que $\Sat{M}{\Gamma}$, mas $\Struct M$ não é uma das
  !!{structure}s que pretendemos.  Isso pode nos levar a adicionar axiomas que
  não são verdadeiros em~$\Struct M$.
\item Se somos bem-sucedidos pelo menos no aspecto de que $\Gamma$ é
  verdadeiro em todas as !!{structure}s pretendidas, então uma sentença~$!A$ é verdadeira em
  todas as !!{structure}s pretendidas sempre que $\Gamma \Entails !A$.  Assim, podemos
  usar ferramentas lógicas (como métodos de !!{derivation}) para mostrar que sentenças são
  verdadeiras em todas as !!{structure}s pretendidas simplesmente mostrando que elas são
  consequência dos axiomas.
\item Às vezes não temos !!{structure}s pretendidas em mente, mas em vez disso
  partimos dos próprios axiomas: começamos com alguns primitivos que
  queremos que satisfaçam certas leis que codificamos em um sistema de axiomas.
  Uma coisa que gostaríamos de verificar imediatamente é que os axiomas
  não se contradizem: se o fizerem, não pode haver conceitos
  que obedeçam a essas leis, e teremos tentado estabelecer uma teoria
  incoerente.  Podemos verificar que isso não acontece encontrando um modelo
  de~$\Gamma$.  E se há modelos da nossa teoria, podemos usar
  métodos lógicos para investigá-los, e também podemos usar
  métodos lógicos para construir modelos.
\item A independência dos axiomas é, igualmente, uma questão
  importante.  Pode acontecer que um dos axiomas seja, na verdade, uma
  consequência dos outros, e seja, portanto, redundante.  Podemos provar que
  um axioma $!A$ em $\Gamma$ é redundante provando $\Gamma \setminus
  \{!A\} \Entails !A$.  Também podemos provar que um axioma não é
  redundante mostrando que $(\Gamma \setminus \{!A\}) \cup \{\lnot
  !A\}$ é satisfatível.  Por exemplo, foi assim que se mostrou que o
  postulado das paralelas é independente dos outros axiomas da geometria.
\item Outra questão importante é a da definibilidade de conceitos
  em uma teoria: A escolha da linguagem determina em que consistem os modelos
  de uma teoria.  Mas nem todo aspecto de uma teoria deve ser
  representado separadamente em seus modelos.  Por exemplo, toda ordenação
  $\le$ determina uma ordenação estrita correspondente~$<$---dada uma, nós
  podemos definir a outra.  Então não é necessário que um modelo de uma
  teoria que envolve tal ordem deva conter \emph{também} a
  ordenação estrita correspondente.  Quando é o caso, em geral,
  de que uma relação pode ser definida em termos de outras?  Quando é
  impossível definir uma relação em termos de outras (e, portanto, deve-se
  adicioná-la aos primitivos da linguagem)?
\end{enumerate}
\end{explain}


\end{document}
